\documentclass[aps,prl,reprint]{revtex4-2}
\usepackage{amsmath}
\usepackage{amssymb}
\usepackage{mathrsfs}
\usepackage{braket}

\begin{document}

\title{On Tensor Representations of Linear Functionals and Bilinear Forms}
\author{A. Scholar}
\affiliation{Department of Mathematics, University of Theory}

\begin{abstract}
This note examines several fundamental questions regarding the representation of linear functionals and bilinear forms in terms of tensor products, covering both finite and infinite-dimensional vector spaces.
\end{abstract}

\maketitle

\section{Linear Functionals via Bilinear Forms}

\subsection{Question}
Can all linear functionals defined on a vector space $V$ be expressed in terms of bilinear functionals with one fixed argument?

\subsection{Answer}
Yes. For any nonzero linear functional $f$ on $V$, pick another nonzero functional $g \in V^*$ and $a \in V$ with $g(a) = 1$. Define the bilinear form:
\[
B(u, v) = g(u)f(v)
\]
Then $B(a, v) = g(a)f(v) = f(v)$ for all $v \in V$.

If $f = 0$, take $B = 0$ and any $a$.

\section{Bilinear Functionals and Tensor Products}

\subsection{Question}
Can all bilinear functionals on a vector space $V$ be expressed as elements of the dual second rank tensor space $V^* \otimes V^*$?

\subsection{Answer}
\begin{itemize}
\item \textbf{Finite-dimensional case:} Yes. If $\dim V = n < \infty$, then:
\[
\text{Bilin}(V, V; \mathbb{F}) \cong V^* \otimes V^*
\]
\item \textbf{Infinite-dimensional case:} No. $V^* \otimes V^*$ contains only finite-rank bilinear functionals. Counterexample: let $\{e_i\}_{i \in I}$ be a basis, define:
\[
B\left(\sum_i a_i e_i, \sum_i b_i e_i\right) = \sum_{i \in I} a_i b_i
\]
This bilinear form is not in $V^* \otimes V^*$.
\end{itemize}

\subsection{Question with $(V \otimes V)^*$}
Can all bilinear functionals be expressed as elements of $(V \otimes V)^*$?

\subsection{Answer}
Yes, by the universal property of the tensor product:
\[
\text{Bilin}(V \times V, \mathbb{F}) \cong (V \otimes V)^*
\]
for any vector space $V$, finite or infinite dimensional.

\section{Inner Products as Tensors}

\subsection{Question}
Can the inner product in a Hilbert space $V$ be considered as a tensor in the dual space $\overline{V}^* \otimes V^*$?

\subsection{Answer}
\begin{itemize}
\item \textbf{Finite-dimensional case:} Yes, the inner product (physics convention: anti-linear in first argument, linear in second) is an element of $\overline{V}^* \otimes V^*$.

\item \textbf{Infinite-dimensional case:} No, not in the algebraic tensor product $\overline{V}^* \otimes V^*$; it belongs to the topological completion.
\end{itemize}

\subsection{Question with $(\overline{V} \otimes V)^*$}
Can the inner product be considered as a tensor in $(\overline{V} \otimes V)^*$?

\subsection{Answer}
Yes. The inner product in physics convention (anti-linear in first argument, linear in second) corresponds exactly to a bilinear map $\overline{V} \times V \to \mathbb{C}$. By the universal property:
\[
\text{Sesquilinear forms} \cong (\overline{V} \otimes V)^*
\]
This holds for any dimension, using the algebraic dual.

\section{Key Definitions}

\subsection{Tensor Product}
The algebraic tensor product $V \otimes W$ is defined by the universal property: bilinear maps $V \times W \to U$ correspond to linear maps $V \otimes W \to U$.

\subsection{Conjugate Space}
For a complex vector space $V$, the conjugate space $\overline{V}$ has the same additive structure but scalar multiplication:
\[
\lambda \cdot_{\overline{V}} v = \overline{\lambda} \cdot_V v
\]

\subsection{Dual Spaces}
\begin{align*}
V^* &= \{ \text{linear functionals } V \to \mathbb{F} \} \\
\overline{V}^* &= \{ \text{anti-linear functionals } V \to \mathbb{C} \}
\end{align*}

\section{Summary}

The representation of linear and bilinear forms via tensor products depends crucially on:
\begin{enumerate}
\item The dimension of the vector space (finite vs. infinite)
\item The specific tensor space used (algebraic tensor product vs. its dual)
\item The convention for inner products (math vs. physics)
\end{enumerate}

The universal property of tensor products provides the most general framework for understanding these relationships.

\end{document} 